%================================================================
\section{Conclusion}\label{sec:conclusion}

A calibrated simulator's answer to a policy question is a function of the
parameter point at which it is evaluated, and the targets that a researcher can
observe rarely pin that point down. In a planted village where the truth is known,
history matching against the status quo and one randomized pilot keeps a large set
of parameter points, and the image of the policy effect over them is wide. The
image, rather than a point inside it, is the report that the targets support.

Against that image we graded intervals computed from one panel, and we kept two
checks apart. The analyst who assumes no interference misses the policy effect at
the real village, and at the calibrator's best-fitting points she misses again and
is overconfident. Analysts who allow for interference through an exposure map and
transport from the pilot cut the bias, at the price of more spread and of intervals
that are conservative relative to it; at the calibrator's points, their one-panel
intervals contain the effect at her point and are still narrower than what the
targets leave open. Containment and interval geometry answer different questions,
and both are needed.

The broader lesson is a reporting rule. A policy number from a calibrated model
should travel with the set of numbers that the same targets cannot rule out, and
an estimate of a network policy from a randomized pilot should travel with the
assumptions that carry it from the pilot to the policy. Both make the answer less
sharp, and both make it easier to believe.
